% Cote 119: offprints, reports, courses, CV and letters (typed and printed part): archive cote 119, site page 303 (PDF page 304).
% typescript. Draft compared: none.
% AI-transcribed by Claude Opus 5.5 (claude-opus-5-5, adaptive thinking, effort high), Message Batches API, batch r119_full (msgbatch_01DwZNm5Y7YsKBWyAN4sFs8w), 2026-11-13; not yet checked by eye.
\scanpage{119/303}
% NOTE: a typed slip is pasted at the head of the page; its text follows
Vous trouverez ci-joint l'exposé
des titres \deleted{et} travaux des candidats
à une direction de recherche

\uline{CURRICULUM VITAE DE ALEXANDRE GROTHENDIECK}

Né le 28 mars 1928 à Berlin, de mère allemande et de père
apatride, émigré de Russie en 1921, mes parents émigrent d'Alle-%
magne en 1933, participent à la révolution espagnole; je les
rejoins en mai 1939. Mes parents sont internés, d'abord mon père
en 1939, puis ma mère en 1940 avec moi. Mon père est déporté du
camp de Vernet en août 1942 pour Auschwitz et est resté disparu;
ma mère meurt en 1957 des suites d'une tuberculose contractée au
camp de concentration. Je reste près de deux ans dans des camps
de concentration français, puis suis recueilli par une maison
d'enfants du ``Secours suisse'' au Chambon-sur-Lignon, où je
termine mes études de lycée en 1945. Etudes de licence (mathé-%
matiques) à Montpellier 1945-48, auditeur libre à l'Ecole Normale
Supérieure à Paris en 1948-49, où je suis le premier séminaire
Cartan sur la théorie des faisceaux, et un cours de Leray \hand{du}
Collège de France sur la théorie de Schauder du degré topologique
dans les espaces localement convexes. De 1949 à 1953 je poursuis
des recherches à Nancy sur les espaces vectoriels topologiques,
comme élève de J. Dieudonné et de L. Schwartz, aboutissant à ma
thèse de doctorat en 1953, sur la théorie des produits tenso-%
riels topologiques et des espaces nucléaires, publiée dans les
``Memoirs of the American Mathematical Society''. Je passe alors
deux ans à l'Université de Saõ Paulo (Brésil), où je continue
et mène à leur aboutissement naturel certaines recherches liées
aux produits tensoriels topologiques [6,7], mais en même temps,
sous l'influence de J. P. Serre, commence à me familiariser avec
des questions de topologie algébrique et d'algèbre homologique.
Ces dernières \unsure{continueront} à m'occuper jusqu'à aujourd'hui, et
sont encore très loin d'être menées à leur terme. Ce sont elles
qui m'occuperont surtout pendant l'année 1955 passée à l'Univer-%
sité du Kansas (USA); j'y développe une théorie commune pour la
théorie de Cartan-Eilenberg des foncteurs dérivés des foncteurs
de modules et la théorie de Leray-Cartan de la cohomologie des
faisceaux [8], et développe des notions de ``cohomologie non commu-%
tative'' dans le contexte des faisceaux et des espaces fibrés à

\mspage{303}
